% Source: source/普通高考/2010/2010安徽文.pdf, page 1, question 13.
% The source PDF is vector line art; the original PNG is retained for reference.
% Nodes: 开始; x=1; x是奇数?; x=x+1; x=x+2; x>8?; 输出x; 结束.
% Directed edges: 开始→x=1→奇偶判定; 是→x=x+1→top loop-back→奇偶判定;
% 否→x=x+2→x>8?; 是→输出x→结束; 否→right-side up-arrow→x=x+1.
\begin{tikzpicture}[
  >=Stealth,
  line width=0.55pt,
  line cap=round,
  line join=round,
  every node/.style={font=\normalsize,anchor=center,align=center,
    execute at begin node={\setlength{\parindent}{0pt}\noindent}},
  flowbox/.style={draw,minimum height=.68cm,inner xsep=4pt,inner ysep=2pt,
    anchor=center,align=center},
  terminal/.style={flowbox,rounded corners=.18cm,minimum width=1.35cm,
    minimum height=.72cm},
  io/.style={flowbox,shape=trapezium,trapezium left angle=72,
    trapezium right angle=108,minimum width=2.50cm,minimum height=.80cm},
  decision/.style={flowbox,shape=diamond,aspect=2.7,minimum width=4.55cm,
    minimum height=1.08cm,inner sep=1pt},
]
  \node[terminal] (start) at (0,11.05) {开始};
  \node[flowbox,minimum width=2.25cm] (init) at (0,9.65) {$x=1$};
  \node[decision,minimum width=4.55cm] (odd) at (0,8.05) {$x\text{ 是奇数?}$};
  \node[flowbox,minimum width=3.20cm] (inc1) at (3.85,8.05) {$x=x+1$};
  \node[flowbox,minimum width=3.20cm] (inc2) at (0,6.35) {$x=x+2$};
  \node[decision,minimum width=3.75cm] (limit) at (0,4.65) {$x>8?$};
  \node[io] (output) at (0,3.05) {输出 $x$};
  \node[terminal] (finish) at (0,1.45) {结束};

  \draw[->] (start.south) -- (init.north);
  % Initial stem joins the shared entry; the single arrow into odd is drawn below.
  \draw (init.south) -- (0,9.00);
  \draw[->] (odd.east) -- node[above=4pt] {是} (inc1.west);
  \draw[->] (odd.south) -- node[right=2pt] {否} (inc2.north);
  \draw[->] (inc2.south) -- (limit.north);
  \draw[->] (limit.south) -- node[right=2pt] {是} (output.north);
  \draw[->] (output.south) -- (finish.north);

  % x>8 的否分支从右侧上行，进入 x=x+1 框底部。
  \coordinate (right-route) at (4.85,4.65);
  \draw (limit.east) -- node[above=2pt] {否} (right-route);
  \draw[->] (right-route) -- (4.85,7.62) -- (inc1.south);

  % 奇数分支更新后由顶部回线返回首个判定框入口。
  \coordinate (top-route) at (6.10,9.00);
  \coordinate (entry) at (0,9.00);
  \draw (inc1.north east) -- (6.10,8.40) -- (top-route);
  \draw[->] (top-route) -- (entry);
  \draw[->] (entry) -- (odd.north);

  \path[use as bounding box] (-2.55,1.05) rectangle (6.50,11.45);
\end{tikzpicture}
